\begin{figure}[t]
\centering
\resizebox{\linewidth}{!}{%
\begin{tikzpicture}[
  font=\sffamily\small,
  root/.style={draw, rounded corners=3pt, fill=blue!10, minimum width=3.4cm, minimum height=1.0cm, align=center, font=\sffamily\bfseries},
  br/.style={draw, rounded corners=3pt, fill=#1, minimum width=3.2cm, minimum height=1.6cm, align=center, text width=3.1cm},
  leaf/.style={align=left, text width=3.2cm, font=\sffamily\scriptsize},
  open/.style={br=orange!15, dashed, thick},
  arr/.style={-{Stealth[length=2mm]}, thick, gray!70},
  node distance=0.35cm and 0.3cm]

\node[root] (r) {Weight-Space\\Programming};

\node[br=green!10, below left=1.2cm and 6.0cm of r] (b1) {\textbf{B1 Reading}\\weights $\to$ description\\{\scriptsize (\nAnyBOne\ studies)}};
\node[br=green!10, right=of b1] (b2) {\textbf{B2 Writing}\\spec $\to$ weights\\{\scriptsize (\nAnyBTwo\ studies)}};
\node[br=green!10, right=of b2] (b3) {\textbf{B3 Repairing}\\weights $\to$ weights\\{\scriptsize (\nAnyBThree\ studies)}};
\node[br=yellow!15, right=of b3] (b4) {\textbf{B4 Feedback \&}\\\textbf{search loops}\\{\scriptsize (\nAnyBFour\ studies)}};
\node[br=yellow!15, right=of b4] (b5) {\textbf{B5 Efficient \&}\\\textbf{non-LLM solvers}\\{\scriptsize (\nAnyBFive\ studies)}};
\node[open, right=of b5] (b6) {\textbf{B6 Self-editing}\\model edits itself\\{\scriptsize (\nAnyBSix\ studies)}};

\foreach \b in {b1,b2,b3,b4,b5,b6} \draw[arr] (r.south) -- (\b.north);

\node[leaf, below=0.25cm of b1] {Original probes and recreations \citep{takehi2026sparks,zen23954378}\\ Prior-vs-reading audits \citep{ax2611_08070,zen24373915}\\ Exactly checkable game policies \citep{ax2611_31979}\\ Orbit audits \citep{ax2611_05798}\\ Clinical and genomics audits \citep{ax2611_05341,ax2611_07847,zen24040591}\\ Trained reader (precursor) \citep{ax2610_07334}
};
\node[leaf, below=0.25cm of b2] {Games and control \citep{ax2611_21449,ax2611_09271,ax2611_16981}\\ Verified policies \citep{ax2611_11604}\\ Physics, quantum, materials \citep{ax2611_12954,ax2611_11537,ax2611_22013}\\ Interpretable transformers \citep{ax2611_18114}\\ Writing ICL \citep{ax2611_26841}
};
\node[leaf, below=0.25cm of b3] {Qwen loop repair and replications \citep{ax2611_09382,ax2611_20953,zen24416417}\\ Robot policies and robot foundation models \citep{ax2611_09376,ax2611_11894,ax2611_31715}\\ VLM, vision, CLIP \citep{ax2611_23615,ax2611_28781,ax2611_25476,ax2611_25257,ax2611_28105}\\ Recommenders and quantized nets \citep{ax2611_29613,ax2611_33031}\\ Transactional edit execution \citep{ax2611_13232}\\ Implanting backdoors \citep{ax2611_17395,ax2611_27920,ax2611_27425}
};
\node[leaf, below=0.25cm of b4] {Experimental feedback \citep{ax2611_07319,ax2611_23718}\\ Federated and training loops \citep{ax2611_08383,ax2611_11005}\\ Null results \citep{ax2611_17789}\\ Governed, reversible loops \citep{ax2611_13236,ax2611_10425}
};
\node[leaf, below=0.25cm of b5] {Ising and LP solvers \citep{ax2611_14210,ax2611_21432}\\ Closed-form and gradient baselines \citep{zen24168084,ax2611_32332,ax2611_33133}\\ Distilled recipes \citep{ax2611_31233,ax2611_29882}\\ Theory \citep{zen23933590,zen24198536}
};
\node[leaf, below=0.25cm of b6] {Rewriting own weights \citep{ax2611_21885}\\ Self-WSP \citep{ax2611_28833}\\ Recursive self-editing \citep{ax2611_29513}\\ Self-upgrading agents \citep{ax2611_31117}
};

\node[draw, rounded corners=3pt, fill=gray!8, font=\sffamily\scriptsize, align=left, text width=4.6cm, above left=0.1cm and 1.2cm of r] (ax) {%
\textbf{Cross-cutting axes}\\
Authorship: I1 direct values, I2 program-mediated, I3 solver-mediated, I4 no LLM\\
Loop: one-shot / closed-loop / search\\
Reference: another model / itself};
\node[draw, rounded corners=3pt, fill=gray!8, font=\sffamily\scriptsize, align=left, text width=4.6cm, above right=0.1cm and 1.2cm of r] (fa) {%
\textbf{Facets}\\
Evaluation (\nFacetEvaluation), verification (\nFacetVerification), safety (\nFacetSafety)\\
Precursors: fast weights, Tracr, hypernetworks, ROME/MEMIT, SEAL};
\end{tikzpicture}}
\caption{Taxonomy of weight-space programming. The six branches describe what the editor does. B1--B3 are the core operations, B4--B5 are how the operations are run and made cheaper, and B6 (dashed) is the open direction connecting WSP to recursive self-improvement (\cref{sec:rsi}). Counts include studies with the branch as primary or secondary code, so a study can be counted in more than one branch. The leaves list representative studies. The full coding is in \cref{tab:papers}.}
\label{fig:taxonomy}
\end{figure}
